\documentclass[10pt,a4paper]{article}

\usepackage[T1]{fontenc}
\usepackage[utf8]{inputenc}
\usepackage{lmodern}
\usepackage{amsmath,amssymb,amsthm}
\usepackage{booktabs,array,tabularx}
\usepackage{geometry}
\usepackage{xcolor}
\usepackage{enumitem}
\usepackage{microtype}
\usepackage{seqsplit}
\usepackage{hyperref}
\usepackage{fancyhdr}

\geometry{margin=1in}
\hypersetup{
  colorlinks=true,
  linkcolor=blue!55!black,
  citecolor=blue!55!black,
  urlcolor=blue!55!black,
  pdftitle={The Eagle Has Wrapped: Practical Public-r Key Recovery against EagleSignV3},
  pdfauthor={Martin Feussner},
  pdfsubject={Cryptanalytic audit report for IACR ePrint 2024/1234},
  pdfkeywords={digital signatures, EagleSignV3, lattices, modulus wrap, key recovery, forgery}
}
\setlist{nosep,leftmargin=*}
\setlength{\parindent}{1em}
\setlength{\parskip}{0.25em}
\setlength{\emergencystretch}{2em}
\pagestyle{fancy}
\fancyhf{}
\fancyhead[L]{The Eagle Has Wrapped}
\fancyhead[R]{15 October 2026}
\fancyfoot[C]{\thepage}
\renewcommand{\headrulewidth}{0.3pt}

\newtheorem{theorem}{Theorem}
\newtheorem{lemma}{Lemma}
\newcommand{\Rq}{R_q}
\newcommand{\HB}{\mathsf{HighBits}_q}
\newcommand{\LB}{\mathsf{LowBits}_q}
\newcommand{\Sign}{\mathsf{Sign}}
\newcommand{\Verify}{\mathsf{Verify}}
\newcommand{\hash}[1]{{\ttfamily\small\seqsplit{#1}}}

\title{\bfseries The Eagle Has Wrapped:\\
Practical Public-$r$ Key Recovery against EagleSignV3}
\author{Martin Feussner\\
\small Selmer Center, University of Bergen\\
\small \texttt{martin.feussner@uib.no}}
\date{15 October 2026}

\begin{document}
\maketitle

\noindent\fcolorbox{black!55}{black!4}{\parbox{0.94\textwidth}{
This audit used the OpenAI model Daybreak Blue at the ultra reasoning level.
OpenAI Codex performed AI-assisted attack-history research, paper review,
cryptanalysis, implementation, experimentation, and technical writing.
OpenAI Codex generated and selected this paper's title and deliberately chose
``The Eagle Has Wrapped'' as wordplay on ``The Eagle Has Landed'' and on the
centered-modulus wrap that exposes the secret-key constraints. Separately
tasked AI agents reviewed bounded parts of the argument and implementation.
The primary full-parameter reproduction, separately tasked hostile adversarial
review, prospectively sealed fresh-context independent reproduction, final
exact-target history review, supplemental public-package replay, and
independent Publication Review all passed. The Publication Reviewer returned
\texttt{PASS / NO VETO}. No independent human verification was performed; it
is optional for this audit.}}

\begin{center}
\fcolorbox{green!45!black}{green!4}{\parbox{0.92\textwidth}{
\textbf{Verified attack.} Independent Publication Review returned
\texttt{PASS / NO VETO}. The package, paper, reproducer, evidence pins, and
release metadata satisfy the audit's publication gates.}}
\end{center}

\begin{abstract}
We describe a practical key-recovery and fresh-message forgery against the
paper-level Ring Level~2 instantiation of EagleSignV3 in IACR ePrint
2024/1234. The recommended row sets
$q=59393$ and $\gamma_2=16384$, although the paper separately requires
$q>4\gamma_2$. Because this inequality is violated, adding the short product
$Dc$ can cross the centered representative boundary at
$\pm(q-1)/2$ while satisfying the signer's printed low-bits test.

For an ordinary returned signature $(r,z)$, the public verifier computes
$v=Tz-Ac=p+Dc\pmod q$. Coordinates of $v$ within $\beta=16$ of a centered
boundary have two possible pre-addition high bits. The public commitment $r$
selects the correct combination. Every selected boundary flip is therefore a
noiseless threshold inequality in one known signed 16-term convolution row of
the ternary secret $D$. Accumulating these inequalities gives a small ternary
integer program. Once $D$ is recovered and public $T$ is a ring unit, the
other signing matrix follows publicly as $G=T^{-1}(A+D)$.

A controlled full-parameter replay retained all 40,000 outputs returned by
the printed signing algorithm. It extracted 21,933 inequalities, including
1,811 wraps, and solved the resulting $21{,}933$-by-$1{,}024$ ternary system
with one SciPy/HiGHS worker in about 68 seconds and 213,404 KiB peak RSS. A
prospectively sealed fresh-key reproduction independently retained 40,000
outputs, extracted 21,982 inequalities including 1,767 wraps, and recovered
the exact signing key in a complete 296.09-second public attack pipeline using at most
221,272 KiB. Both candidates satisfied every inequality and the public
relation before hidden-key comparison, then produced accepted signatures on
unqueried messages with changed-message and altered-signature controls
rejecting. The runs used, respectively, the Appendix's 32-byte and the main
text's 64-byte completion of $r$. Each is one exact-row demonstration; neither
estimates a minimum query count or an all-key success probability. No exact
EagleSignV3 source or known-answer vector was located, so the claim is
mathematical and paper-level rather than an implementation-conformance result.
\end{abstract}

\section{Introduction}

EagleSignV3 is a Fiat--Shamir-with-aborts lattice signature designed as a
successor to EagleSign V1 and V2~\cite{target}. Its public key is
$(\rho,T)$, where $\rho$ expands to a public matrix $A$ and
\[
  T=(A+D)G^{-1}\pmod q
\]
for short secret matrices $D$ and $G$. The signature changes the predecessors'
response to $z=y+Gc$ and applies rejection tests to $z$ and the low bits of
$p+Dc$, where $p=Ty$. The paper presents this as the change that restores the
zero-knowledge property after the attacks of Pulles and
Tibouchi~\cite{pullestibouchi}.

The Ring Level~2 parameter row does not satisfy a condition stated in the same
paper. It chooses $q=59393$ and $\gamma_2=2^{14}$, so
\[
  q=59393 < 65536 = 4\gamma_2.
\]
This matters because EagleSignV3 first reduces coefficients to centered
representatives modulo $q$ and then decomposes them modulo
$2\gamma_2=32768$. A short addition near $\pm(q-1)/2$ can wrap through the
centered modulus boundary and change the high part by two. The low part still
lies well inside the signer's acceptance interval.

The hash value $r$ commits to the high bits before $Dc$ is added. The public
signature and public key reveal the high bits after that addition. Near a
boundary, testing the few possible wrap patterns against $r$ recovers exact
comparison outcomes for coefficients of $Dc$. Because $c$ has Hamming weight
16 and $D$ is ternary, each outcome is a sparse integer inequality in the
1,024 secret coefficients. Standard mixed-integer feasibility then recovers
$D$ and, for unit $T$, the equivalent signing key $G$.

We make no claim that boundary analysis, sparse linear inequalities, or
mixed-integer solving is a new cryptanalytic technique. The result is their
application to the exact EagleSignV3 Ring Level~2 equations and parameters,
followed by a full key-recovery and ordinary-verifier fresh-message forgery.
The final history result is a bounded search statement, not an absolute
firstness or attack-technique novelty claim.

\paragraph{Contributions.}
We isolate the exact centered-wrap condition, prove that public $r$ labels the
resulting sparse inequalities, give a public-only recovery procedure, and
report two complete one-key full-parameter executions under distinct
$r$-length completions. We explicitly separate the paper-level mathematical
result from byte-level completion choices and from the behavior of any
unavailable author implementation.

\section{Exact target and claim boundary}
\label{sec:target}

\subsection{Pinned paper}

The target is Hounkpevi, Djimnaibeye, Seck, and Sow's
\emph{EagleSignV3: A new secure variant of EagleSign signature over lattices},
IACR ePrint 2024/1234, revised 6 August 2024~\cite{target}. The pinned 35-page
PDF has SHA-256
\begin{center}
\hash{9b346bf2234299eadc72460a47c20dd71e6131fa19b3d21abdaaea01a8f061a2}.
\end{center}

No exact public EagleSignV3 implementation or known-answer vector was located.
The paper says that an implementation can be deduced from EagleSignV2, but V2
is a predecessor with different signing equations and an explicit boundary
guard. It is not a conformance oracle for this target. All executable claims
in this report therefore concern an explicit reconstruction of the printed V3
mathematical scheme.

\subsection{Frozen Ring Level~2 row}

The attacked row is the ring variant at the paper's NIST security level~2:
\begin{table}[ht]
\centering
\small
\begin{tabular}{@{}ll@{}}
\toprule
Quantity & Frozen value \\
\midrule
Ring & $\mathbb Z_q[X]/(X^{1024}+1)$ \\
$(k,l,n)$ & $(1,1,1024)$ \\
$q$ & $59393$ \\
Challenge weight $\tau$ & $16$ \\
Secret coefficient bound $\eta$ & $1$ \\
$\beta=l\eta\tau$ & $16$ \\
$(\gamma_1,\gamma_2)$ & $(16384,16384)$ \\
Printed classical strong-unforgeability estimate & $117$ bits \\
Printed classical long-term key-recovery estimate & $172$ bits \\
\bottomrule
\end{tabular}
\caption{The exact paper-level parameter row used by the attack.}
\label{tab:parameters}
\end{table}

The paper defines centered representatives modulo odd $q$ in
$[-(q-1)/2,(q-1)/2]$, and centered representatives modulo even
$2\gamma_2$ in $(-\gamma_2,\gamma_2]$. We use those conventions literally.
Decomposing unsigned residues directly would describe a different scheme.

The paper's hash and byte descriptions conflict in several places. In
particular, Section~3.3 gives a 512-bit range for hash $G$, making $r$ 64
bytes, while Appendix~A.4 says 256 bits, making $r$ 32 bytes. The main signing
algorithm and appendix also disagree on the length of $\lambda$. The primary
replay used the Appendix's 32-byte $r$ completion. The prospectively sealed
independent reproduction and the proposed public reproducer use the main
text's 64-byte $r$ completion and Appendix~A.4's 48-byte $\lambda$. Both
profiles fix all remaining byte choices before execution. The wrap argument
depends only on the printed ring equations, centered decomposition, challenge
distribution, and public commitment to the pre-addition high bits; its success
under both $r$ lengths is therefore paper-level evidence across the conflict,
not evidence for a canonical wire format.

\subsection{Output semantics}

The claim consumes every value returned after Algorithm~4's two printed
rejection checks. In the primary run, 1,767 of 40,000 returned values---about
4.4\%---were rejected when subsequently submitted to the ordinary verifier.
That mismatch is another consequence of the broken wrap invariant. The
returned values still expose public $(r,z)$ and therefore valid wrap labels
for the recovery system.

An implementation that adds an extra ordinary self-verification test and
resamples failed outputs defines a different signing procedure. We do not
claim that the present 40,000-output experiment transfers unchanged to that
modified procedure. This boundary is essential to the result.

\section{EagleSignV3 equations}

Let $\Rq=\mathbb Z_q[X]/(X^n+1)$. At the attacked ring row all matrices are
single ring elements. Key generation samples ternary $D,G\in\Rq$, with $G$
invertible, expands $\rho$ to uniform $A\in\Rq$, and publishes
\[
  T=(A+D)G^{-1}\pmod q.
\]

To sign a message $M$, the signer computes a message digest $\mu$, samples
$y\in S_{\gamma_1}$, and forms
\[
  p=Ty,\qquad r=G_{\rm hash}(\HB(p)),\qquad
  c=H(\mu,r),\qquad z=y+Gc.
\]
The challenge $c$ has exactly $\tau=16$ coefficients in $\{-1,+1\}$ and all
other coefficients zero. The signer returns $(r,z)$ after checking
\[
  \lVert z\rVert_\infty < \gamma_1-\beta,
  \qquad
  \lVert\LB(p+Dc)\rVert_\infty < \gamma_2-\beta.
\]

The verifier recomputes $c=H(\mu,r)$ and
\[
  v=Tz-Ac.
\]
Using $TG=A+D$ and $z=y+Gc$ gives the central public identity
\begin{equation}
  v=Ty+(TG-A)c=p+Dc\pmod q.          \label{eq:identity}
\end{equation}
It accepts when $z$ is short and
$c=H(\mu,G_{\rm hash}(\HB(v)))$.

The paper separately instructs designers to choose
$q>4\max(\gamma_1,\gamma_2)$. That condition would place enough room between
the decomposition scale and the centered $q$ boundary for the intended
high-bit stability argument. Table~\ref{tab:parameters} violates it.

\section{Public wrap labels}
\label{sec:labels}

Write
\[
  h=(q-1)/2=29696,\qquad B=2\gamma_2=32768,
  \qquad s=(Dc)_j.
\]
Because $D$ is ternary and $c$ has weight 16, $|s|\leq16=\beta$. A coordinate
of the public $v$ can have crossed the centered boundary only when its
representative lies within 16 units of $-h$ or $+h$.

\begin{lemma}[Lower boundary]
Suppose $v_j=-h+a$ with $0\leq a<\beta$. The hidden pre-addition coefficient
$p_j$ lies on the positive side of the centered boundary exactly when
$s>a$.
\end{lemma}
\begin{proof}
Equation~\eqref{eq:identity} gives $p_j=v_j-s\pmod q$. Its centered
representative wraps from the displayed negative value to the positive side
exactly when $v_j-s<-h$. Substituting $v_j=-h+a$ yields $s>a$.
\end{proof}

\begin{lemma}[Upper boundary]
Suppose $v_j=h-a$ with $0\leq a<\beta$. The hidden pre-addition coefficient
$p_j$ lies on the negative side of the centered boundary exactly when
$s<-a$.
\end{lemma}
\begin{proof}
Here $p_j=v_j-s$ crosses above $h$ exactly when
$h-a-s>h$, or $s<-a$.
\end{proof}

For a lower-edge coordinate, $\HB(v_j)=-1$ while the wrapped alternative for
$\HB(p_j)$ is $+1$. For an upper edge the public value is $+1$ and the wrapped
alternative is $-1$. If one signature contains $K$ edge coordinates, the
attacker enumerates their $2^K$ possible toggles, hashes each complete
high-bit vector, and compares with public $r$. A unique match identifies every
wrap bit. In the primary run, $K$ ranged from zero to five; all 40,000 outputs
had a unique match, and only 69,252 candidate vectors were tested in total.

\section{From labels to the signing key}

Let the nonzero challenge coefficients be $(c_b)_{b\in S}$ with
$|S|=16$. Negacyclic multiplication gives
\[
  (Dc)_j=\sum_{b\in S}\epsilon_{j,b}c_bD_{j-b\bmod n},
\]
where $\epsilon_{j,b}=+1$ before the $X^n=-1$ fold and $-1$ after it. Every
coefficient is therefore a known signed sum of 16 unknowns
$D_i\in\{-1,0,1\}$.

The recovered lower-edge label supplies either
\[
  (Dc)_j\geq a+1 \quad\text{or}\quad (Dc)_j\leq a,
\]
and the upper-edge label supplies either
\[
  (Dc)_j\leq-a-1 \quad\text{or}\quad (Dc)_j\geq-a.
\]
These are exact integer inequalities, with no statistical noise term. Pooling
them over returned signatures gives a ternary feasibility problem in 1,024
variables.

After solving for a candidate $D$, the attacker performs two public checks.
First, every collected inequality must hold exactly. Second, when public $T$
is a unit in $\Rq$, compute
\begin{equation}
  G=T^{-1}(A+D)\pmod q.              \label{eq:gderive}
\end{equation}
The centered coefficients of $G$ must be ternary, $G$ must be invertible, and
$TG=A+D$ must hold. These checks select a signing-equivalent pair using public
data alone; comparison with the generated hidden key is unnecessary for the
attack and is postponed until after the forgery gate.

For the paper-level Ring Level~2 distribution, the estimated probability that
$T$ is a unit is
\[
  (1-1/q)^n\approx0.9829.
\]
The direct bridge in Equation~\eqref{eq:gderive} is explicitly conditional on
this public property. A nonunit $T$ does not invalidate the inequality
recovery, but deriving a complete signer would need another method and is not
claimed here.

\begin{theorem}[Conditional signing compromise]
For the paper-level Ring Level~2 equations, suppose returned signatures give a
ternary $D$ satisfying all public wrap inequalities, public $T$ is a unit, and
Equation~\eqref{eq:gderive} yields a ternary invertible $G$ satisfying
$TG=A+D$. Then $(D,G)$ is an equivalent signing key and can run the printed
signing equations on a fresh message.
\end{theorem}
\begin{proof}
The public checks establish the exact key-generation relation for the same
$A,T$. Substituting the recovered pair into signing yields
$Tz-Ac=Ty+(TG-A)c=Ty+Dc$, exactly the correctness identity used by ordinary
verification. Rejection sampling can be repeated until the printed bounds and
ordinary verification accept. No hidden value from the original key is
needed.
\end{proof}

\section{Full-parameter evidence}
\label{sec:evidence}

\subsection{Primary reproduction: Appendix 32-byte \texorpdfstring{$r$}{r}
completion}

The controlled primary replay fixed one Ring Level~2 key and retained exactly
40,000 outputs of the printed signing algorithm. Table~\ref{tab:primary}
summarizes the stable result.

\begin{table}[ht]
\centering
\small
\begin{tabular}{@{}lr@{}}
\toprule
Measurement & Result \\
\midrule
Printed signer outputs & 40,000 \\
Public edge inequalities & 21,933 \\
Centered-boundary wraps & 1,811 \\
Ordinary-verifier-invalid returned outputs & 1,767 (4.4175\%) \\
Hash candidates tested & 69,252 \\
Ternary variables / nonzeros & 1,024 / 350,928 \\
SciPy/HiGHS branch-and-bound nodes & 1 \\
SciPy/HiGHS solver / command wall time & 65.67 s / 68.08 s \\
SciPy/HiGHS peak RSS / workers & 213,404 KiB / 1 \\
Fresh-message ordinary verification & accepted \\
Changed-message / altered-$r$ controls & rejected / rejected \\
\bottomrule
\end{tabular}
\caption{Primary paper-level Ring Level~2 reproduction.}
\label{tab:primary}
\end{table}

Collection took 598.77 seconds. The regenerated constraint dataset has
SHA-256
\begin{center}
\hash{a15776674731b97cca2c469ab4d43512c966d6d4636e5e59344ffd8d38ffb97c}.
\end{center}
SciPy/HiGHS returned a ternary candidate satisfying all 21,933 inequalities. Before
any hidden-key comparison, the implementation derived ternary invertible $G$,
checked $TG=A+D$, and signed the never-queried message
\begin{center}
\texttt{eaglesignv3-highs-continuation-fresh-message-never-queried}.
\end{center}
The ordinary verifier accepted that signature. Reusing it on a changed
message failed, and flipping one bit of $r$ while retaining $z$ also failed.
Only after these public checks passed did a post-hoc comparison confirm exact
recovery of the generated $D$ and $G$.

The fresh-forgery and recovered-key artifacts have SHA-256 values
\begin{center}
\hash{1e90009b4eb0d9e9cb0103c6015ee7b3a76bb47a6a4334b0ac703fa79eeb5822},\\
\hash{e0e0fa88ff3bef951687344d17731018495cb958bf68e50d65f7cff3d0bf4eb1}.
\end{center}

The fixed count 40,000 is one demonstrated working point. We did not optimize
the number of outputs and make no minimum-query claim. The experiment used one
fixed key and therefore does not measure a distribution-wide success rate.

\subsection{Hostile adversarial review}

A separately tasked reviewer worked with the explicit objective of disproving
the attack and returned
\texttt{PASS\_WITH\_LIMITATIONS / NO\_VETO}. The review covered the exact
Ring Level~2 arithmetic, both strict wrap inequalities, public-$r$ label
selection, the sparse negacyclic convolution rows, ternary solver bounds,
public derivation of $G$, fresh-message signing, ordinary verification, and
the primary endpoint. It found no hidden-key input or ground-truth label in
candidate selection. The separate verifier shares the audited public model
code, which remains a disclosed evidentiary limitation.

The review preserved the material boundaries stated here: paper-level
Ring Level~2 only; one fixed primary key; direct signing conditional on unit
$T$; 40,000 as a working point rather than a minimum; use of every output
returned by printed Algorithm~4; explicit completion of conflicting byte
details; and no author-implementation claim. Its handoff, full review,
verdict, and sealed artifact manifest have SHA-256 values
\begin{center}
\hash{e8c51f43aff9c1a57bd44a619bb718715ee2d340c9e29c01b46228a524f881cb},\\
\hash{af6fec885f924f5df83e202fec9b05106c897fe53d5edb7a7a1f1e2218e218dc},\\
\hash{e08c94ef6da1dcbc3a66f46cd231054a345b056ab30b524e6f09fa36638bb420},\\
\hash{4e1bcf3d6607514189b4b1362805e8b361c5433cc774b54af405547d049daad7}.
\end{center}
This gate authorizes the sealed independent reproduction; it does not itself
count or publish the attack.

\subsection{Fresh-context independent reproduction: main-text 64-byte
\texorpdfstring{$r$}{r} completion}

A fresh-context implementation and run were prospectively sealed before the
target seed was created. The seal covered the hypothesis, fixed completion
profile, exact commands, C++20 source, SciPy/HiGHS solver wrapper, and all six
executables. The fresh 32-byte kernel-CSPRNG seed was generated only after that
seal. Key attempt zero had unit public $T$ and was the only generated key.

The independent completion follows Section~3.3's 512-bit definition of hash
$G$, so $r$ is 64 bytes. It follows Appendix~A.4 for 48-byte $\lambda$ and
32-byte $\rho$ and $tr$, uses 48-byte $\mu$, and samples uniformly from the
inclusive set $[-\gamma_1,\gamma_1]$. This differs deliberately from the
primary replay's Appendix-based 32-byte $r$ completion while preserving the
same printed Ring Level~2 equations and attack mechanism.

The run retained every one of 40,000 outputs returned after Algorithm~4's two
printed rejection tests. Of those, 38,273 passed and 1,727 failed the ordinary
verifier; no verifier filter was applied. Public edge enumeration produced
21,982 exact ternary inequalities, including 1,767 edge flips, from 69,301
candidate high-bit vectors. The largest single transcript had six edge
coordinates. SciPy~1.18.1 \texttt{optimize.milp}, using bundled
HiGHS~1.12.0, solved the 1,024-variable, 351,712-nonzero system in 141.57
solver seconds at one node with zero gap and zero checked row violation.

Before hidden-key access, the implementation checked ternarity, every
inequality, unit $G$, and $TG=A+D$, then signed the never-queried message
\begin{center}
\texttt{EagleSignV3 independent fresh forgery 2026-10-15 / ESV3-IR-001}.
\end{center}
The unchanged ordinary verifier accepted the untouched forgery and rejected
both a changed message and an altered signature. Only afterward did a separate
process compare against ground truth and confirm exact $D$ and $G$ matches.

The oracle took 124.51 seconds, constraint extraction 12.83 seconds, the
SciPy/HiGHS command 144.56 seconds, recovery and fresh signing 12.82 seconds,
and public verification 1.37 seconds. The complete public pipeline therefore
took 296.09 seconds, excluding post-hoc comparison, with one worker, one-CPU
solver affinity, a 4 GiB virtual-memory ceiling, and 221,272 KiB peak RSS.

The prospective seal, run artifact manifest, solver report, recovered key,
fresh forgery, and terminal handoff have SHA-256 values
\begin{center}
\hash{2e8144384c31f59a3dd4d2d196f2d51da885c8bd8fde25e0d7713440e2b67934},\\
\hash{19da75184dde5c2697d8f796bdd3e33a0b09c7700318151ff4c0a9e42613de8a},\\
\hash{4e66d7a7c57bdc073088e5c63320c31fb00b0b4b448a5b2514dc639ac89dedc9},\\
\hash{4dc18591edf542ccd5ea1f5ef13cb01f00973694b9a1e62510bac4a193acc737},\\
\hash{aad6f1610b96b5d5b5d10e961d334b1a240fc7f7d4e5066ae7a42fc7316345c5},\\
\hash{bf51fc87cc1c7f2060c1e30caea121d4f57183e2496424a3646f00672e0473b0}.
\end{center}

\subsection{Publication Review and supplemental package replay}

Independent Publication Review returned \texttt{PASS / NO VETO}. As
supplemental release-readiness validation, the complete public-package path
was also replayed on a new fresh unit-$T$ key. It retained all 40,000 printed
signer outputs: 38,219 passed ordinary verification and 1,781 failed it, with
no verifier filtering. The replay produced 22,083 exact inequalities including
1,821 wraps, recovered exact $D$ and $G$, checked the public relation, and
produced an accepted fresh-message forgery. The unchanged verifier rejected
the changed message and altered signature. The one-worker solve took exactly
83.31292366300477 solver-wrapper seconds and 84.10 command-wall seconds, with
243,116 KiB peak RSS.

This package replay is supplemental validation of the release path. It does
not replace the prospectively sealed independent reproduction, enlarge the
claim beyond the primary and independent two-key trial scope, or change the
reported two trials and two successes. The replay review, machine-readable
report, retained output manifest, and review artifact manifest have SHA-256
values
\begin{center}
\hash{ee6fefb703b8be1c50f5d2804c57f5a62bc1f7f377f1bf8ca0f7aef3769b7761},\\
\hash{b1549fd2456b8f68d8393ddca1cf23e1ecd2cb43fc12cddb4337832796583e9a},\\
\hash{b969b537a02e87439a2d07d25a44592e23fb5a187526bbcec9ff6ba8f4427fad},\\
\hash{8437db885f97a5021353cff69bef8b1cf7604ddbc200feb458196938da369837}.
\end{center}
The replay input was the pre-finalization source candidate, pinned by root
manifest
\hash{9c8fecf3ba1ea16df8ebd60c52d117a014b0939ac1bb9c85a9fea511a25c93d5}
and nested reproducer manifest
\hash{d5b48642fe37e2a043af4b6e9a79219452e9a73d4583c31d4c8a297579e249ac}.
These are old source-candidate pins, not hashes of this final package.

\section{Reproduction architecture}

The public reproducer is an independently written C++20 paper-level
implementation. It uses separate executables for key generation, oracle
collection, public recovery, ordinary verification, and post-hoc ground-truth
comparison. The recovery executable receives only the public key and the
returned $(M,r,z)$ records. It writes an LP-format integer program. A sealed
Python wrapper parses that file into a sparse matrix and invokes SciPy~1.18.1
\texttt{optimize.milp} with bundled HiGHS~1.12.0, integral ternary bounds,
one thread, and one-CPU affinity. The wrapper independently validates every
inequality before emitting its candidate. Recovery then derives $G$ publicly
and creates the fresh signature. The verifier separately checks freshness,
the public relation, ordinary acceptance, changed-message rejection, and an
altered-signature rejection before ground truth is consulted.

The package offers a structural smoke test and a full fresh-key run:
\begin{verbatim}
./reproduce.sh --smoke
./reproduce.sh --full
\end{verbatim}
Full mode uses one worker, a 4 GiB virtual-memory ceiling, 40,000 returned
outputs, and the sealed 21,600-second SciPy/HiGHS time limit. Generated secret
material and the large transcript remain in a temporary directory and are
deleted after the final post-hoc comparison.

Because the paper has no exact V3 source and contains conflicting byte
details, the reproducer records its choices: a 64-byte $r$ following the main
text, 48-byte $\lambda$, 32-byte $\rho$ and $tr$, 48-byte $\mu$, uniform
inclusive $[-\gamma_1,\gamma_1]$ sampling, canonical little-endian public
coefficients, and a documented high-bit encoding. These choices preserve the
printed equations and distributions but do not define a canonical V3 wire
format.

\section{Attack history and attribution}
\label{sec:history}

The final post-reproduction history review returned \texttt{PASS}, with
moderate confidence and no history veto. Its approved bounded statement is:
\begin{quote}
No prior public substantive attack against EagleSignV3 (IACR ePrint
2024/1234) was located in the searches performed as of 15 October 2026.
\end{quote}
This statement cannot exclude private, removed, unpublished, or unindexed
work, and search and citation indexes can lag.

The direct-lineage V1/V2 cryptanalysis is documented in IACR ePrint
2024/1137~\cite{pullestibouchi}.
Pulles and Tibouchi previously broke EagleSign V1 and V2 using response statistics and ideal structure. Those schemes use different response relations, and EagleSign V2 includes an explicit centered-modulus-boundary rejection that is absent from the EagleSignV3 specification studied here. We located no earlier public disclosure applying the centered-modulus-wrap, public-commitment-label route to EagleSignV3. This is a bounded search statement; we make no absolute firstness or attack-technique novelty claim.

Prabowo and Tan analyze generic Schnorr-type signatures with aborts and give a
statistical attack on the original EagleSign~\cite{prabowotan}. Their work
provides context for response statistics but does not attack the exact V3
scheme or analyze the centered-wrap public-$r$ route. Pellet-Mary and Seck
study generic properties of HighBits and LowBits and their
applications~\cite{pelletmaryseck}. The retrieved abstract, keywords, visible references,
and indexed search surfaces disclosed no EagleSign attack; we treat that work
as adjacent technique context rather than an exact-target attack.

The canonical history handoff, review handoff, machine-readable verdict,
query ledger, source ledger, and review artifact manifest have SHA-256 values
\begin{center}
\hash{2a1cbf638453500b5749cbc040ece0d03abb453e0bb95f7ccd7f4d4f2df82faa},\\
\hash{53b5c8f3ca4726764b9e0db6df5f6c458f25d988b693968044245a01bfe31d62},\\
\hash{b28391fda40777c64919a61035a33f8135eccb60a2f635cc7e78d793e2093236},\\
\hash{43b6eb77e399f8119fc88b162db57ff80d2d8dbf4a5ab127bfcfe4af700256c2},\\
\hash{f71cd94c1e1e6a9a3af690e39c63074e61aaf98aea7e1bd16705951b5c2f2f48},\\
\hash{d50ad9e9bebf5a1855c89f960592e131f19fce2b87d8bb5e4514f567d217454d}.
\end{center}

The attack technique itself is not claimed as original. The contribution is
the exact application of standard boundary reasoning and integer feasibility
to the published EagleSignV3 Ring Level~2 row, with a complete signing-key
bridge and fresh-message forgery.

\section{Limitations and repair boundary}

The following qualifications are part of the claim:
\begin{itemize}
  \item The result covers only the exact paper-level Ring Level~2 row in
        Table~\ref{tab:parameters}. It does not claim the module rows, Level~5
        rows, or any unpublished implementation.
  \item No exact V3 source or known-answer vector was located. Conflicting
        hash lengths, seed lengths, and encodings required explicit
        completions; there is no author-source or wire-format conformance
        claim.
  \item The direct signing-key bridge applies when public $T$ is a unit,
        estimated at about 98.29\% for this paper-level row and detectable
        from the public key.
  \item The primary and independent evidence are each one-key executions.
        They use different keys, implementations, and $r$-length completions,
        but together remain a two-key record rather than a many-key
        success-rate study.
  \item Forty thousand outputs were sufficient in both reported runs; they
        are not asserted to be necessary or minimal.
  \item The attack uses every output of the printed signer, including the
        roughly 4.4\% that ordinary verification rejects. An added
        self-verification-and-resampling loop is outside the claim.
  \item The final history PASS has moderate confidence and is a dated bounded
        search result. It cannot establish absolute priority or exclude
        private, deleted, unpublished, or unindexed work.
\end{itemize}

The paper's own condition $q>4\max(\gamma_1,\gamma_2)$ excludes the reported
parameter mismatch. Enforcing that condition in a revised row, or restoring a
correct centered-boundary guard before returning signatures, are natural
starting points for repair. Either change alters the analyzed scheme or its
parameters and requires fresh correctness, rejection-rate, security, and
implementation analysis. Adding only a final self-verification loop blocks
the precise output stream used here, but is likewise a modified signer and is
not by itself a security proof.

\section{Conclusion}

At EagleSignV3 Ring Level~2, the printed modulus is too small for the paper's
own decomposition condition. Public $r$ resolves the resulting centered-wrap
choices and converts ordinary signer outputs into exact sparse inequalities in
the ternary secret $D$. One full-parameter primary run solved 21,933 such
inequalities under the Appendix's 32-byte $r$ completion. A prospectively
sealed fresh-key reproduction independently solved 21,982 inequalities under
the main text's 64-byte $r$ completion. Each derived $G$ from the public
relation and produced an accepted fresh-message signature with no hidden-key
input to recovery. The result is a practical paper-level signing compromise on
two demonstrated keys, subject to the explicit unit-$T$, output-semantics, and
specification-completion limits above. Independent Publication Review returned
\texttt{PASS / NO VETO}, and the supplemental public-package replay confirmed
the complete release path without changing the two-key claim scope.

\section*{Acknowledgements}
The computations were performed on the Norwegian Research and Education Cloud (NREC), using resources provided by the University of Bergen and the University of Oslo.

\begin{thebibliography}{9}

\bibitem{target}
A. C. Hounkpevi, S. Djimnaibeye, M. Seck, and D. Sow,
``EagleSignV3: A new secure variant of EagleSign signature over lattices,''
IACR Cryptology ePrint Archive, Report 2024/1234, revised 6 August 2024.
\url{https://eprint.iacr.org/2024/1234}

\bibitem{pullestibouchi}
L. N. Pulles and M. Tibouchi,
``Cryptanalysis of EagleSign,''
IACR Cryptology ePrint Archive, Report 2024/1137, 2024.
\url{https://eprint.iacr.org/2024/1137}

\bibitem{prabowotan}
T. F. Prabowo and C. H. Tan,
``Concrete Analysis of Schnorr-type Signatures with Aborts,''
IACR Cryptology ePrint Archive, Report 2024/1256, 2024.
\url{https://eprint.iacr.org/2024/1256}

\bibitem{pelletmaryseck}
A. Pellet-Mary and M. Seck,
``On the Properties of HighBits and LowBits Functions and their
Applications,''
IACR Cryptology ePrint Archive, Report 2026/895, 2026.
\url{https://eprint.iacr.org/2026/895}

\bibitem{highs}
Q. Huangfu and J. A. J. Hall,
``Parallelizing the dual revised simplex method,''
\emph{Mathematical Programming Computation}, vol.~10, pp.~119--142, 2018.

\end{thebibliography}

\end{document}
